\documentclass[a4paper]{article}
\usepackage[fontset=fandol]{ctex}
\usepackage[margin=1.65cm]{geometry}
\usepackage{amsmath,amssymb}
\usepackage{graphicx}
\usepackage{hyperref}
\usepackage{longtable,booktabs}
\usepackage{xcolor}
\usepackage{enumitem}
\usepackage{listings}
\lstset{basicstyle=\ttfamily\small,breaklines=true,columns=fullflexible,keepspaces=true}
\hypersetup{colorlinks=true,linkcolor=blue,urlcolor=blue}
\setlist[itemize]{leftmargin=1.25em,itemsep=0.18em}
\setlength{\parskip}{0.45em}
\setlength{\parindent}{2em}
\sloppy
\title{第 15 讲：Mid/Post-training：SFT 与 RLHF}
\author{CS336 Spring 2026 public videos and official course materials}
\date{}
\begin{document}
\maketitle
\tableofcontents
\newpage
\section{本章导读}
本章讨论 \textbf{Mid/Post-training：SFT 与 RLHF}。CS336 的第一层目标不是把现成大模型当作黑箱使用，而是把 language model 从输入文本、训练目标、模型结构、数据、系统和评估这些部件重新拆开。只有拆开之后，读者才能判断一个设计选择究竟改变了什么。
设想 base model 会续写文本，却不会稳定遵循 instruction。Post-training（后训练）要把语言模型变成 assistant：SFT 提供示范行为，RLHF/DPO 用偏好约束输出风格和安全边界。
本章对应 CS336 Spring 2026 第 15 个公开视频，并结合课程页提供的可解析资料。视频和课程页给出事实边界；正文把这些材料改写为可自学的教材章节。
阅读本章时，先理解问题为什么存在，再看定义、公式和实现形态，随后通过例题和练习检查自己是否能独立复现这条 reasoning chain。
\begin{figure}[h]
\centering
\includegraphics[width=0.58\linewidth]{source_anchor.jpg}
\caption{第 15 讲的课程图像锚点。}
\end{figure}
\subsection{本章知识路线}
\begin{itemize}
\item 本章首先说明 Mid/Post-training：SFT 与 RLHF 为什么是 language modeling pipeline 中不可跳过的一环。
\item 其次定义本章反复使用的术语，并说明这些术语对应的计算对象或系统对象。
\item 随后用公式和伪代码表达主要机制，让概念能够落到可计算的变量上。
\item 最后给出例题、风险讨论和练习，便于读者检查自己是否真正掌握了机制。
\end{itemize}
\subsection{结构图}
下面的图用于帮助读者定位本章的核心结构。先看概念之间的依赖，再回到正文中的公式、伪代码和例题。
\begin{figure}[h]
\centering
\includegraphics[width=0.72\linewidth]{pdf_pages/page-02.jpg}
\caption{官方材料图页：Mid/Post-training：SFT 与 RLHF 的课程页/PDF/script 图像锚点。}
\end{figure}
\begin{figure}[h]
\centering

\includegraphics[width=0.72\linewidth]{pdf_pages/page-01.jpg}
\caption{官方材料图页：Mid/Post-training：SFT 与 RLHF 的课程页/PDF/script 图像锚点。}
\end{figure}
\section{本章主线}
本章可以沿着三条线索阅读：课程为什么把这个主题放在这里，它改变了哪些变量，以及怎样用小实验检查这些变量。
\subsection{本章要解决的核心问题}
本讲从 pretraining 转向 mid/post-training。SFT 用示范回答塑造 instruction-following，RLHF 用偏好数据进一步控制输出。
这一点对应的学习动作是：找出对象，写出变量，说明变量如何被测量。若输入规模、硬件、数据分布或评估协议改变，结论也必须重新检查。
因此，Mid/Post-training：SFT 与 RLHF 不应被读成若干名词的集合，而应被读成一组可检验的机制。
\subsection{从机制到资源账本}
post-training 更 artisanal：数据稀缺、公开细节少、标注规范和 reward model 都会强烈影响模型行为。
这一点对应的学习动作是：找出对象，写出变量，说明变量如何被测量。若输入规模、硬件、数据分布或评估协议改变，结论也必须重新检查。
因此，Mid/Post-training：SFT 与 RLHF 不应被读成若干名词的集合，而应被读成一组可检验的机制。
\subsection{从课堂讲解到可复现实验}
DPO/PPO/RLHF 等方法的共同主题是让模型偏向更符合人类偏好的回答，同时用 KL 或 reference policy 控制偏离。
这一点对应的学习动作是：找出对象，写出变量，说明变量如何被测量。若输入规模、硬件、数据分布或评估协议改变，结论也必须重新检查。
因此，Mid/Post-training：SFT 与 RLHF 不应被读成若干名词的集合，而应被读成一组可检验的机制。
\section{术语、符号与问题边界}
本章的术语横跨模型、数据和系统。读者需要先知道每个词指向的对象：有些词描述数学目标，有些词描述张量形状，有些词描述硬件瓶颈，还有一些词描述评估协议。术语边界不清，后面的公式和实现就会变成空洞的缩写。
\begin{longtable}{p{0.08\linewidth}p{0.42\linewidth}p{0.42\linewidth}}
\toprule
\textbf{\#} & \textbf{术语} & \textbf{阅读要求}\\
\midrule
1 & post-training（后训练） & 需要能说明它指向的对象、参与的公式或代码位置，以及一个典型失败情形。\\
2 & supervised fine-tuning (SFT) & 需要能说明它指向的对象、参与的公式或代码位置，以及一个典型失败情形。\\
3 & RLHF & 需要能说明它指向的对象、参与的公式或代码位置，以及一个典型失败情形。\\
4 & reward model（奖励模型） & 需要能说明它指向的对象、参与的公式或代码位置，以及一个典型失败情形。\\
5 & DPO & 需要能说明它指向的对象、参与的公式或代码位置，以及一个典型失败情形。\\
6 & PPO & 需要能说明它指向的对象、参与的公式或代码位置，以及一个典型失败情形。\\
\bottomrule
\end{longtable}
这些术语之间有依赖关系。例如 tokenization 改变 sequence length，sequence length 又改变 attention FLOPs、KV cache 和 evaluation token budget；parallelism 改变显存占用，也改变通信瓶颈和故障恢复方式。
\subsection{关键词}
下面的关键词在课程材料中反复出现。它们不代替定义，只提示读者本章哪些对象需要在后文持续跟踪。
\begin{itemize}
\item data, rlhf, reward, many, instruction, safety, monopsony, market, feedback, approach, part, high, examples, good
\end{itemize}
\section{从 pretraining 到 assistant 行为}
现在进入 从 pretraining 到 assistant 行为。本节关心的不是名称本身，而是它把 Mid/Post-training：SFT 与 RLHF 中的哪些对象变成可计算、可测量或可优化的量。
\subsection{为什么要讨论 从 pretraining 到 assistant 行为}
视频开场说 post-training 比 pretraining 更 messy、更 artisanal；它把基础 LM 转成可交互助手。
pretraining 学 next-token distribution，SFT 学示范回答格式和任务行为。
mid-training/post-training 的核心是改变模型可见行为，而不是重新学习全部世界知识。
从 pretraining 到 assistant 行为 的作用是把 Mid/Post-training：SFT 与 RLHF 中的一个抽象问题变成可计算对象：哪些变量进入公式，哪些状态进入代码，哪些条件会让结果失效。
\subsection{从 pretraining 到 assistant 行为 的机制}
视频开场说 post-training 比 pretraining 更 messy、更 artisanal；它把基础 LM 转成可交互助手。 这给出本节的对象和边界。读者需要把其中的名词对应到张量、数据记录、训练状态或系统状态，而不是停留在缩写层面。
pretraining 学 next-token distribution，SFT 学示范回答格式和任务行为。 因而同一个方法在小规模示例、单机实验和集群训练中的意义并不相同。比较方法时，问题规模、硬件条件、数据分布和评估口径必须同时给出。
mid-training/post-training 的核心是改变模型可见行为，而不是重新学习全部世界知识。 这使本节具有可检验性：如果一个结论不能在公式、伪代码、profile trace、benchmark 或 ablation 中表现出来，它就还只是直觉。
因此，从 pretraining 到 assistant 行为 应当被理解为一组可以实现和测试的假设。CS336 反复强调 from scratch，正是因为只有能被复现的机制，才可能被稳定改进。
\subsection{从 pretraining 到 assistant 行为 的数学表达}
\[\mathcal{L}_{\text{SFT}}(\theta)=-\sum_{(x,y)}\sum_t \log p_\theta(y_t\mid x,y_{<t})\]
符号说明：\(x\) 是 instruction/context，\(y\) 是人工或高质量示范回答；SFT 仍是 token-level 监督学习。
这个关系式给出了本节最小的数学骨架。读者应检查每个变量的单位、取值范围和可观测性；如果某个量只能间接估计，后续实验就必须说明 proxy 是什么。
例如，validation loss 常被用作模型质量的 proxy，tokens/s 用作吞吐 proxy，Jaccard 或 MinHash 用作近重复 proxy，reward model score 用作偏好 proxy。proxy 不是目标本身；它只在评估协议清楚时才可引用。
\subsection{从 从 pretraining 到 assistant 行为 到程序}
本节机制可以写成下面的伪代码。它保留课程材料中的计算顺序和状态变化，但省略了工程代码中的错误处理、并行细节和框架样板。
\begin{lstlisting}[language=Python]
for instruction, response in sft_data:
    loss = cross_entropy(model(instruction, response), response_tokens)
    update(loss)
\end{lstlisting}
读这段伪代码时，重点不是语法，而是状态如何变化：输入是什么，中间变量是什么，主要计算或数据移动发生在哪里，哪些边界条件会破坏结果。
\subsection{边界条件与常见误解}
\begin{itemize}
\item SFT 数据质量比数量更重要，坏示范会被直接模仿。
\item SFT 可能降低 base model 的某些原始能力或校准。
\item 不要把小实验中的局部结论自动外推到 frontier-scale；scale、数据分布和硬件拓扑都会改变结论。
\item 不要把 benchmark 或 profile 的单个数字当作机制解释；必须同时记录实验设置、输入形状、版本和评估口径。
\end{itemize}
\subsection{与本章其他概念的关系}
从 pretraining 到 assistant 行为 与本章其他部分的关系是互补的：术语表给出对象名称，公式给出最小数学结构，伪代码给出执行顺序，例题则把这些内容压缩成一个可检查的计算。
对这一类主题，最常见的误解是只列方法名。更有用的做法是比较每个方法改变的成本项、依赖的假设和可能的失败模式。
\subsection{怎样检验 从 pretraining 到 assistant 行为}
检验 从 pretraining 到 assistant 行为 的第一步是确定测量对象。这里的测量对象通常不是一个抽象名词，而是与 post-training、SFT、instruction、assistant、demonstration 有关的张量、样本、请求、奖励、通信操作或评估记录。如果对象没有被具体化，后面的实验就只能得到描述性观察，不能得到可复现结论。
第二步是构造最小输入。最小输入应该足够小，使读者能手算或打印关键中间量；同时又要保留本节机制。对于 从 pretraining 到 assistant 行为，最小输入不追求逼真规模，而追求暴露因果关系：改变一个变量时，哪些输出、成本或错误会随之改变。
第三步是写出资源账本。账本可以是 FLOPs、bytes moved、显存峰值、通信量、token 数、样本数、reward 分布或 benchmark 分数。符号说明：\textbackslash{}(x\textbackslash{}) 是 instruction/context，\textbackslash{}(y\textbackslash{}) 是人工或高质量示范回答；SFT 仍是 token-level 监督学习。 这类说明应被转化为单位明确的数量关系；如果数量只能间接观测，就必须说明使用了什么 proxy。
第四步是设置对照。一个好对照只改变一个因素，例如 batch size、sequence length、vocabulary size、attention heads、filter threshold、learning rate、decoding 参数或 verifier。如果多个因素同时改变，实验最多说明系统表现不同，不能说明是哪一个机制在起作用。
第五步是观察失败条件。教材中的成功路径往往最容易记住，但工程中真正决定方法边界的是失败案例。本节至少要记住这些限制：SFT 数据质量比数量更重要，坏示范会被直接模仿；SFT 可能降低 base model 的某些原始能力或校准。复习时可以把这些限制改写成反例测试。
最后，把实验结论放回 Mid/Post-training：SFT 与 RLHF 的整体结构中。从 pretraining 到 assistant 行为 不是孤立技巧：它会影响训练目标、数据选择、系统成本或评估解释中的至少一项。能说清楚这种影响，才说明读者已经从名词记忆进入机制理解。
\subsection{常见混淆：从 pretraining 到 assistant 行为}
post-training 在本节中应被看作一个技术对象，而不是一个孤立名词。需要判断它指向模型、数据、系统还是评估，并说明它在“从 pretraining 到 assistant 行为”中改变了哪个变量或约束。
supervised fine-tuning (SFT) 用示范回答塑造 assistant 行为。在“从 pretraining 到 assistant 行为”中，这个概念用于解释 Mid/Post-training：SFT 与 RLHF 的一个具体机制。
instruction 在本节中应被看作一个技术对象，而不是一个孤立名词。需要判断它指向模型、数据、系统还是评估，并说明它在“从 pretraining 到 assistant 行为”中改变了哪个变量或约束。
assistant 在本节中应被看作一个技术对象，而不是一个孤立名词。需要判断它指向模型、数据、系统还是评估，并说明它在“从 pretraining 到 assistant 行为”中改变了哪个变量或约束。
demonstration 在本节中应被看作一个技术对象，而不是一个孤立名词。需要判断它指向模型、数据、系统还是评估，并说明它在“从 pretraining 到 assistant 行为”中改变了哪个变量或约束。
这些概念的共同点是：只有进入公式、代码或评估协议之后才有明确含义。单独背诵名称不会帮助判断一个训练 run、推理系统或数据 pipeline 是否正确。
\subsection{从 pretraining 到 assistant 行为 的小结}
\begin{itemize}
\item 从 pretraining 到 assistant 行为 的任何性能结论都必须同时报告问题规模和评估口径，否则不可比较。
\item 如果一个方法声称降低主要成本，它通常会把成本转移到另一个变量，例如显存、通信、数据质量、训练稳定性或评估复杂度。
\item 公式和伪代码提供推理骨架；真实系统仍需要 profiler、ablation 和 held-out evaluation。
\end{itemize}
\section{Reward modeling 与 RLHF}
现在进入 Reward modeling 与 RLHF。本节关心的不是名称本身，而是它把 Mid/Post-training：SFT 与 RLHF 中的哪些对象变成可计算、可测量或可优化的量。
\subsection{为什么要讨论 Reward modeling 与 RLHF}
RLHF 用人类偏好训练 reward model，再用 RL 优化策略模型。
pairwise preference 更容易标注开放式回答质量，但 reward model 只是偏好的近似。
PPO 等 RL 方法需要 KL penalty，防止模型远离 reference policy 而产生 reward hacking。
后训练部分的核心问题是行为控制。模型已经学到语言分布，但还需要通过示范、偏好或可验证奖励，把输出推向可用、可评估、可约束的行为。
\subsection{Reward modeling 与 RLHF 的机制}
RLHF 用人类偏好训练 reward model，再用 RL 优化策略模型。 这给出本节的对象和边界。读者需要把其中的名词对应到张量、数据记录、训练状态或系统状态，而不是停留在缩写层面。
pairwise preference 更容易标注开放式回答质量，但 reward model 只是偏好的近似。 因而同一个方法在小规模示例、单机实验和集群训练中的意义并不相同。比较方法时，问题规模、硬件条件、数据分布和评估口径必须同时给出。
PPO 等 RL 方法需要 KL penalty，防止模型远离 reference policy 而产生 reward hacking。 这使本节具有可检验性：如果一个结论不能在公式、伪代码、profile trace、benchmark 或 ablation 中表现出来，它就还只是直觉。
因此，Reward modeling 与 RLHF 应当被理解为一组可以实现和测试的假设。CS336 反复强调 from scratch，正是因为只有能被复现的机制，才可能被稳定改进。
\subsection{Reward modeling 与 RLHF 的目标函数}
\[\mathcal{L}_{\text{RM}}=-\log\sigma(r_\phi(x,y^+)-r_\phi(x,y^-))\]
符号说明：\(y^+\) 是偏好回答，\(y^-\) 是较差回答，\(r_\phi\) 是 reward model。
这个关系式给出了本节最小的数学骨架。读者应检查每个变量的单位、取值范围和可观测性；如果某个量只能间接估计，后续实验就必须说明 proxy 是什么。
例如，validation loss 常被用作模型质量的 proxy，tokens/s 用作吞吐 proxy，Jaccard 或 MinHash 用作近重复 proxy，reward model score 用作偏好 proxy。proxy 不是目标本身；它只在评估协议清楚时才可引用。
\subsection{Reward modeling 与 RLHF 的训练流程}
本节机制可以写成下面的伪代码。它保留课程材料中的计算顺序和状态变化，但省略了工程代码中的错误处理、并行细节和框架样板。
\begin{lstlisting}[language=Python]
reward = reward_model(prompt, response)
kl = kl_divergence(policy, reference)
objective = reward - beta * kl
policy.update(objective)
\end{lstlisting}
读这段伪代码时，重点不是语法，而是状态如何变化：输入是什么，中间变量是什么，主要计算或数据移动发生在哪里，哪些边界条件会破坏结果。
\subsection{边界条件与常见误解}
\begin{itemize}
\item reward model 会被策略优化 exploited。
\item 偏好数据反映标注规范和人群，不是客观真理。
\item 不要把小实验中的局部结论自动外推到 frontier-scale；scale、数据分布和硬件拓扑都会改变结论。
\item 不要把 benchmark 或 profile 的单个数字当作机制解释；必须同时记录实验设置、输入形状、版本和评估口径。
\end{itemize}
\subsection{与本章其他概念的关系}
Reward modeling 与 RLHF 与本章其他部分的关系是互补的：术语表给出对象名称，公式给出最小数学结构，伪代码给出执行顺序，例题则把这些内容压缩成一个可检查的计算。
对这一类主题，最常见的误解是只列方法名。更有用的做法是比较每个方法改变的成本项、依赖的假设和可能的失败模式。
\subsection{怎样检验 Reward modeling 与 RLHF}
检验 Reward modeling 与 RLHF 的第一步是确定测量对象。这里的测量对象通常不是一个抽象名词，而是与 RLHF、reward model、preference、PPO、KL 有关的张量、样本、请求、奖励、通信操作或评估记录。如果对象没有被具体化，后面的实验就只能得到描述性观察，不能得到可复现结论。
第二步是构造最小输入。最小输入应该足够小，使读者能手算或打印关键中间量；同时又要保留本节机制。对于 Reward modeling 与 RLHF，最小输入不追求逼真规模，而追求暴露因果关系：改变一个变量时，哪些输出、成本或错误会随之改变。
第三步是写出资源账本。账本可以是 FLOPs、bytes moved、显存峰值、通信量、token 数、样本数、reward 分布或 benchmark 分数。符号说明：\textbackslash{}(y\textasciicircum{}+\textbackslash{}) 是偏好回答，\textbackslash{}(y\textasciicircum{}-\textbackslash{}) 是较差回答，\textbackslash{}(r\_\textbackslash{}phi\textbackslash{}) 是 reward model。 这类说明应被转化为单位明确的数量关系；如果数量只能间接观测，就必须说明使用了什么 proxy。
第四步是设置对照。一个好对照只改变一个因素，例如 batch size、sequence length、vocabulary size、attention heads、filter threshold、learning rate、decoding 参数或 verifier。如果多个因素同时改变，实验最多说明系统表现不同，不能说明是哪一个机制在起作用。
第五步是观察失败条件。教材中的成功路径往往最容易记住，但工程中真正决定方法边界的是失败案例。本节至少要记住这些限制：reward model 会被策略优化 exploited；偏好数据反映标注规范和人群，不是客观真理。复习时可以把这些限制改写成反例测试。
最后，把实验结论放回 Mid/Post-training：SFT 与 RLHF 的整体结构中。Reward modeling 与 RLHF 不是孤立技巧：它会影响训练目标、数据选择、系统成本或评估解释中的至少一项。能说清楚这种影响，才说明读者已经从名词记忆进入机制理解。
\subsection{常见混淆：Reward modeling 与 RLHF}
RLHF 用人类偏好 reward 优化模型行为。在“Reward modeling 与 RLHF”中，这个概念用于解释 Mid/Post-training：SFT 与 RLHF 的一个具体机制。
reward model 在本节中应被看作一个技术对象，而不是一个孤立名词。需要判断它指向模型、数据、系统还是评估，并说明它在“Reward modeling 与 RLHF”中改变了哪个变量或约束。
preference 在本节中应被看作一个技术对象，而不是一个孤立名词。需要判断它指向模型、数据、系统还是评估，并说明它在“Reward modeling 与 RLHF”中改变了哪个变量或约束。
PPO 在本节中应被看作一个技术对象，而不是一个孤立名词。需要判断它指向模型、数据、系统还是评估，并说明它在“Reward modeling 与 RLHF”中改变了哪个变量或约束。
KL 在本节中应被看作一个技术对象，而不是一个孤立名词。需要判断它指向模型、数据、系统还是评估，并说明它在“Reward modeling 与 RLHF”中改变了哪个变量或约束。
这些概念的共同点是：只有进入公式、代码或评估协议之后才有明确含义。单独背诵名称不会帮助判断一个训练 run、推理系统或数据 pipeline 是否正确。
\subsection{Reward modeling 与 RLHF 的小结}
\begin{itemize}
\item Reward modeling 与 RLHF 的任何性能结论都必须同时报告问题规模和评估口径，否则不可比较。
\item 如果一个方法声称降低主要成本，它通常会把成本转移到另一个变量，例如显存、通信、数据质量、训练稳定性或评估复杂度。
\item 公式和伪代码提供推理骨架；真实系统仍需要 profiler、ablation 和 held-out evaluation。
\end{itemize}
\section{DPO 与直接偏好优化}
现在进入 DPO 与直接偏好优化。本节关心的不是名称本身，而是它把 Mid/Post-training：SFT 与 RLHF 中的哪些对象变成可计算、可测量或可优化的量。
\subsection{为什么要讨论 DPO 与直接偏好优化}
DPO 类方法把偏好优化写成监督式目标，避免显式训练 RL loop。
它依然依赖偏好数据和 reference policy，只是把 reward 隐式化。
post-training 方法选择应看数据类型、稳定性、compute、目标行为和安全约束。
后训练部分的核心问题是行为控制。模型已经学到语言分布，但还需要通过示范、偏好或可验证奖励，把输出推向可用、可评估、可约束的行为。
\subsection{DPO 与直接偏好优化 的机制}
DPO 类方法把偏好优化写成监督式目标，避免显式训练 RL loop。 这给出本节的对象和边界。读者需要把其中的名词对应到张量、数据记录、训练状态或系统状态，而不是停留在缩写层面。
它依然依赖偏好数据和 reference policy，只是把 reward 隐式化。 因而同一个方法在小规模示例、单机实验和集群训练中的意义并不相同。比较方法时，问题规模、硬件条件、数据分布和评估口径必须同时给出。
post-training 方法选择应看数据类型、稳定性、compute、目标行为和安全约束。 这使本节具有可检验性：如果一个结论不能在公式、伪代码、profile trace、benchmark 或 ablation 中表现出来，它就还只是直觉。
因此，DPO 与直接偏好优化 应当被理解为一组可以实现和测试的假设。CS336 反复强调 from scratch，正是因为只有能被复现的机制，才可能被稳定改进。
\subsection{DPO 与直接偏好优化 的目标函数}
\[\mathcal{L}_{\text{DPO}}=-\log\sigma\left(\beta\log\frac{\pi_\theta(y^+\mid x)}{\pi_{\text{ref}}(y^+\mid x)}-\beta\log\frac{\pi_\theta(y^-\mid x)}{\pi_{\text{ref}}(y^-\mid x)}\right)\]
符号说明：\(\pi_\theta\) 是待训练策略，\(\pi_{\text{ref}}\) 是参考模型，\(\beta\) 控制偏离参考模型的强度。
这个关系式给出了本节最小的数学骨架。读者应检查每个变量的单位、取值范围和可观测性；如果某个量只能间接估计，后续实验就必须说明 proxy 是什么。
例如，validation loss 常被用作模型质量的 proxy，tokens/s 用作吞吐 proxy，Jaccard 或 MinHash 用作近重复 proxy，reward model score 用作偏好 proxy。proxy 不是目标本身；它只在评估协议清楚时才可引用。
\subsection{DPO 与直接偏好优化 的训练流程}
本节机制可以写成下面的伪代码。它保留课程材料中的计算顺序和状态变化，但省略了工程代码中的错误处理、并行细节和框架样板。
\begin{lstlisting}[language=Python]
chosen_logp = policy.logprob(prompt, chosen)
rejected_logp = policy.logprob(prompt, rejected)
loss = dpo_loss(chosen_logp, rejected_logp, reference_logps)
\end{lstlisting}
读这段伪代码时，重点不是语法，而是状态如何变化：输入是什么，中间变量是什么，主要计算或数据移动发生在哪里，哪些边界条件会破坏结果。
\subsection{边界条件与常见误解}
\begin{itemize}
\item DPO 不是万能替代 RL；它受偏好数据覆盖范围限制。
\item post-training 改善 helpfulness 可能牺牲 calibration 或 diversity。
\item 不要把小实验中的局部结论自动外推到 frontier-scale；scale、数据分布和硬件拓扑都会改变结论。
\item 不要把 benchmark 或 profile 的单个数字当作机制解释；必须同时记录实验设置、输入形状、版本和评估口径。
\end{itemize}
\subsection{与本章其他概念的关系}
DPO 与直接偏好优化 与本章其他部分的关系是互补的：术语表给出对象名称，公式给出最小数学结构，伪代码给出执行顺序，例题则把这些内容压缩成一个可检查的计算。
对这一类主题，最常见的误解是只列方法名。更有用的做法是比较每个方法改变的成本项、依赖的假设和可能的失败模式。
\subsection{怎样检验 DPO 与直接偏好优化}
检验 DPO 与直接偏好优化 的第一步是确定测量对象。这里的测量对象通常不是一个抽象名词，而是与 DPO、preference optimization、reference、beta、alignment 有关的张量、样本、请求、奖励、通信操作或评估记录。如果对象没有被具体化，后面的实验就只能得到描述性观察，不能得到可复现结论。
第二步是构造最小输入。最小输入应该足够小，使读者能手算或打印关键中间量；同时又要保留本节机制。对于 DPO 与直接偏好优化，最小输入不追求逼真规模，而追求暴露因果关系：改变一个变量时，哪些输出、成本或错误会随之改变。
第三步是写出资源账本。账本可以是 FLOPs、bytes moved、显存峰值、通信量、token 数、样本数、reward 分布或 benchmark 分数。符号说明：\textbackslash{}(\textbackslash{}pi\_ heta\textbackslash{}) 是待训练策略，\textbackslash{}(\textbackslash{}pi\_\{ ext\{ref\}\}\textbackslash{}) 是参考模型，\textbackslash{}(\textbackslash{}beta\textbackslash{}) 控制偏离参考模型的强度。 这类说明应被转化为单位明确的数量关系；如果数量只能间接观测，就必须说明使用了什么 proxy。
第四步是设置对照。一个好对照只改变一个因素，例如 batch size、sequence length、vocabulary size、attention heads、filter threshold、learning rate、decoding 参数或 verifier。如果多个因素同时改变，实验最多说明系统表现不同，不能说明是哪一个机制在起作用。
第五步是观察失败条件。教材中的成功路径往往最容易记住，但工程中真正决定方法边界的是失败案例。本节至少要记住这些限制：DPO 不是万能替代 RL；它受偏好数据覆盖范围限制；post-training 改善 helpfulness 可能牺牲 calibration 或 diversity。复习时可以把这些限制改写成反例测试。
最后，把实验结论放回 Mid/Post-training：SFT 与 RLHF 的整体结构中。DPO 与直接偏好优化 不是孤立技巧：它会影响训练目标、数据选择、系统成本或评估解释中的至少一项。能说清楚这种影响，才说明读者已经从名词记忆进入机制理解。
\subsection{常见混淆：DPO 与直接偏好优化}
DPO 把偏好优化写成直接监督式目标。在“DPO 与直接偏好优化”中，这个概念用于解释 Mid/Post-training：SFT 与 RLHF 的一个具体机制。
preference optimization 在本节中应被看作一个技术对象，而不是一个孤立名词。需要判断它指向模型、数据、系统还是评估，并说明它在“DPO 与直接偏好优化”中改变了哪个变量或约束。
reference 在本节中应被看作一个技术对象，而不是一个孤立名词。需要判断它指向模型、数据、系统还是评估，并说明它在“DPO 与直接偏好优化”中改变了哪个变量或约束。
beta 在本节中应被看作一个技术对象，而不是一个孤立名词。需要判断它指向模型、数据、系统还是评估，并说明它在“DPO 与直接偏好优化”中改变了哪个变量或约束。
alignment 在本节中应被看作一个技术对象，而不是一个孤立名词。需要判断它指向模型、数据、系统还是评估，并说明它在“DPO 与直接偏好优化”中改变了哪个变量或约束。
这些概念的共同点是：只有进入公式、代码或评估协议之后才有明确含义。单独背诵名称不会帮助判断一个训练 run、推理系统或数据 pipeline 是否正确。
\subsection{DPO 与直接偏好优化 的小结}
\begin{itemize}
\item DPO 与直接偏好优化 的任何性能结论都必须同时报告问题规模和评估口径，否则不可比较。
\item 如果一个方法声称降低主要成本，它通常会把成本转移到另一个变量，例如显存、通信、数据质量、训练稳定性或评估复杂度。
\item 公式和伪代码提供推理骨架；真实系统仍需要 profiler、ablation 和 held-out evaluation。
\end{itemize}
\section{综合讨论：从局部机制到完整系统}
本章的主题是 Mid/Post-training：SFT 与 RLHF，但它不应被读成几个相邻知识点的拼接。从 pretraining 到 assistant 行为、Reward modeling 与 RLHF、DPO 与直接偏好优化 共同回答同一个问题：语言模型系统中哪些对象可以被定义，哪些成本可以被估算，哪些结论必须通过实验或评估来确认。
这些对象包括 post-training（后训练）、supervised fine-tuning (SFT)、RLHF、reward model（奖励模型）、DPO、PPO。读者可以把它们看成一组接口：有的接口连接文本和 token，有的连接模型结构和张量计算，有的连接数据样本和训练目标，有的连接离线评估和在线服务。接口一旦改变，后面的成本、错误类型和优化空间也会随之改变。
\subsection{从 pretraining 到 assistant 行为 在系统中的位置}
从 pretraining 到 assistant 行为 首先提供一个局部解释：视频开场说 post-training 比 pretraining 更 messy、更 artisanal；它把基础 LM 转成可交互助手。 这类解释的价值在于把直觉压缩成可操作对象。一旦对象明确，读者就可以追问它的输入、输出、中间状态和失败条件，而不是只记住方法名。
其次，从 pretraining 到 assistant 行为 会改变至少一个资源或评估变量。它可能改变 sequence length、activation size、memory bandwidth、communication volume、data mixture、reward distribution 或 benchmark protocol。这些变量在小例子里可以手算，在真实训练和服务系统里则需要 profiling、logging 和 held-out evaluation。
最后，从 pretraining 到 assistant 行为 不能脱离边界条件理解。一个关键 caveat 是：SFT 数据质量比数量更重要，坏示范会被直接模仿。。把 caveat 写出来不是为了削弱结论，而是为了说明结论在哪些输入、硬件、数据和评估条件下才成立。
\subsection{Reward modeling 与 RLHF 在系统中的位置}
Reward modeling 与 RLHF 首先提供一个局部解释：RLHF 用人类偏好训练 reward model，再用 RL 优化策略模型。 这类解释的价值在于把直觉压缩成可操作对象。一旦对象明确，读者就可以追问它的输入、输出、中间状态和失败条件，而不是只记住方法名。
其次，Reward modeling 与 RLHF 会改变至少一个资源或评估变量。它可能改变 sequence length、activation size、memory bandwidth、communication volume、data mixture、reward distribution 或 benchmark protocol。这些变量在小例子里可以手算，在真实训练和服务系统里则需要 profiling、logging 和 held-out evaluation。
最后，Reward modeling 与 RLHF 不能脱离边界条件理解。一个关键 caveat 是：reward model 会被策略优化 exploited。。把 caveat 写出来不是为了削弱结论，而是为了说明结论在哪些输入、硬件、数据和评估条件下才成立。
\subsection{DPO 与直接偏好优化 在系统中的位置}
DPO 与直接偏好优化 首先提供一个局部解释：DPO 类方法把偏好优化写成监督式目标，避免显式训练 RL loop。 这类解释的价值在于把直觉压缩成可操作对象。一旦对象明确，读者就可以追问它的输入、输出、中间状态和失败条件，而不是只记住方法名。
其次，DPO 与直接偏好优化 会改变至少一个资源或评估变量。它可能改变 sequence length、activation size、memory bandwidth、communication volume、data mixture、reward distribution 或 benchmark protocol。这些变量在小例子里可以手算，在真实训练和服务系统里则需要 profiling、logging 和 held-out evaluation。
最后，DPO 与直接偏好优化 不能脱离边界条件理解。一个关键 caveat 是：DPO 不是万能替代 RL；它受偏好数据覆盖范围限制。。把 caveat 写出来不是为了削弱结论，而是为了说明结论在哪些输入、硬件、数据和评估条件下才成立。
\subsection{把本章知识用于读论文和复现实验}
读相关论文时，可以把本章变成一个检查表。第一，论文是否清楚定义了输入规模、模型规模、数据来源和评估协议；第二，论文声称的改进落在哪个变量上；第三，作者是否报告了会让结论失效的设置；第四，实验是否能分离模型结构、数据质量、优化设置和系统实现的影响。
复现实验时，也应避免直接追求大规模。先用最小程序复现一个公式、一个伪代码分支或一个 profile 现象，再扩大输入规模。这样做的原因很简单：如果小规模程序无法解释中间量，大规模训练只会把错误隐藏在日志和随机波动里。
把这些检查合在一起，本章给出的不是一个固定答案，而是一种阅读 CS336 的方法：每个模型机制都要落到变量，每个变量都要落到测量，每个测量都要说明协议，每个协议都要承认边界。
因此，本章中的任何一句结论都可以被改写成一个可引用的完整句式：在给定数据、模型、硬件和评估协议下，某个机制改变了某个变量，并通过某个观测指标表现出来。这样的句子比单独背诵术语更长，但它包含了自学、复现和引用时真正需要的信息。
把本章用于自学时，可以按三遍阅读。第一遍只追踪对象和定义，确认每个术语到底指向文本、token、张量、样本、请求还是评估记录。第二遍追踪公式和伪代码，确认每一步改变了哪些状态。第三遍追踪实验和 caveat，确认哪些结论只在特定规模、数据或硬件条件下成立。
把本章用于复习时，则应反过来操作：先从一个失败案例或异常指标出发，再回到相应机制。loss 曲线异常可能指向数据、优化或模型容量；吞吐异常可能指向 kernel、通信或调度；benchmark 异常可能指向 contamination、prompt protocol 或采样设置。能够从异常反推机制，是掌握本章的实用标准。
\section{例题：把概念变成可计算检查}
\subsection{DPO 目标如何偏向 chosen response}
同一 prompt 有 chosen 和 rejected 两个回答，以及 reference model 的 log probability。
\begin{itemize}
\item DPO 比较 policy 相对 reference 对 chosen/rejected 的偏好差。
\item 如果 chosen 的相对 log probability 提高，loss 下降。
\item 参数 `β` 控制偏离 reference 的力度。
\end{itemize}
\textbf{结论。} 偏好优化仍然依赖数据覆盖和 reference policy，不能把目标函数当成自动 alignment。
这个例题的目的不是替代完整工程实现，而是给出一种最小推理形式：把问题转成变量，写出近似公式或伪代码，再说明近似何时失效。
\section{实验设计：从小例子到可复现结论}
学习 Mid/Post-training：SFT 与 RLHF 不能停在概念层。一个教材化结论至少需要四件事：可构造的输入、可观测的中间量、可比较的指标，以及清楚的失败条件。下面的实验设计不是要求训练大模型，而是要求读者用小规模程序复现本章机制。
\subsection{最小输入}
先构造只触发本章核心机制的小输入。例如 tokenizer 章节可以用几个包含 rare words 和 Unicode 字符的字符串；GPU 章节可以用一个 elementwise kernel 和一个 GEMM；data 章节可以用十几篇近重复文档。小输入的价值在于它让错误可见。
\subsection{资源账本}
对同一输入写下 FLOPs、bytes moved、显存峰值、通信量或样本数量的估算。估算不需要一开始就精确，但必须说明单位和近似假设。随后用 profiler、benchmark、validation loss 或人工检查去验证估算是否同量级。
\subsection{单变量消融}
一次只改变一个因素，例如 vocabulary size、head 数、batch size、filter threshold、reward scale、decoding temperature 或 GPU 数量。若多个变量同时变化，实验只能说明系统变了，不能说明机制是什么。
\subsection{失败条件}
最后要刻意触发失败：极端 context length、错误 dtype、过小 batch、污染 benchmark、低质量数据或不可靠 verifier。失败条件常常比成功曲线更能说明一个方法的边界。
\subsection{记录与报告}
实验记录应能让另一个读者复现同一结论。最少要写清楚输入构造、代码版本、随机种子、硬件、batch shape、数据版本、评价指标和异常处理方式。报告结论时不要只写“更快”或“更好”，而要说明快在哪里、好在哪个指标上，以及是否牺牲了内存、稳定性、数据质量或可解释性。
常见报告错误是把局部现象写成普遍规律。例如一次小模型实验里的 loss 改善，不等于更大模型或不同数据分布上也会改善；一次 kernel benchmark 的吞吐提升，不等于端到端 serving latency 一定降低。教材中的实验报告应始终把结论限定在可复现条件内。
\section{课程资料与回看路径}
教材正文已经把主要概念、公式和实现逻辑组织成连续叙述；本节只提供回到课程材料的阅读路径。读者复习时可以先完成前文练习，再按下面的时间段或资料组核对细节。
\subsection{视频回看路径}
\begin{longtable}{p{0.20\linewidth}p{0.46\linewidth}p{0.26\linewidth}}
\toprule
\textbf{时间} & \textbf{关键词} & \textbf{复习任务}\\
\midrule
00:00:05-00:07:01 & going, very, data, think, data., it's, pre-training & 回看定义、例子、推导或 caveat 在该段中的位置。\\
00:20:45-00:27:19 & think, knowledge, don't, it's, very, might, your & 回看定义、例子、推导或 caveat 在该段中的位置。\\
00:40:35-00:46:32 & good, going, that's, talk, pre-training, your, very & 回看定义、例子、推导或 caveat 在该段中的位置。\\
00:53:43-01:00:14 & that's, annotators, think, there's, it's, people, quite & 回看定义、例子、推导或 caveat 在该段中的位置。\\
01:13:09-01:19:48 & very, think, rlhf, take, going, good, step & 回看定义、例子、推导或 caveat 在该段中的位置。\\
\bottomrule
\end{longtable}
这些时间段用于定位视频中的讲解顺序。严肃引用时，应回到视频片段确认讲者表述、板书或幻灯片内容。
回看视频时，不必逐字抄录字幕。更有效的方法是把每个时间段判定为定义、例子、推导、代码解释或 caveat，再把它放回本章对应小节。
\subsection{课程资料阅读路径}
\begin{longtable}{p{0.14\linewidth}p{0.48\linewidth}p{0.30\linewidth}}
\toprule
\textbf{资料组} & \textbf{标题/页块} & \textbf{关键词}\\
\midrule
official\_01 & Lecture 15 “A F T E R P R E T R A I N I N G ” ( M I D / P O S T T R A I N I N G ) ... & data, information, annotation, source, part, cs336, class\\
official\_03 & Style variations in data and models Models vary a lot in response length.. [Wang+ ... & knowledge, monopsony, market, style, relevant, teaching, extraction\\
official\_05 & RLHF Data What types of pairwise feedback. How do we get (good) pairwise feedback? & rlhf, data, pairwise, feedback, standard, types, good\\
official\_06 & Modern worker distribution Caveat: many platforms, covers only one (Outlier, ScaleAI) & rlhf, data, distribution, annotators, feedback, crowdsourcing, pairwise\\
official\_08 & DPO in LLaMA (and other open models) DPO + Expert iteration for post-training & rlhf, variants, many, watch, pref, lots, overoptimization\\
\bottomrule
\end{longtable}
课程资料通常给出更稳定的标题、公式、图或代码结构；视频则补充动机和口头解释。两者合起来构成本章的事实边界。
阅读资料时，可以把每个资料组拆成三个对象：一个定义，一个公式或伪代码动作，一个边界条件。这样比逐字翻译页面或脚本更容易形成可用知识。
\section{实践说明、历史位置与风险}
Mid/Post-training：SFT 与 RLHF 在 CS336 的课程结构中连接基础机制和规模化问题。基础机制解释方法为什么成立；规模化问题检验它在更大模型、更长上下文、更复杂数据或真实 serving workload 下是否仍成立。
历史位置不等于模型年表。更重要的是识别问题如何迁移：早期方法解决小规模建模问题，现代方法常常解决同一问题在大规模数据、GPU 集群、长上下文或人类偏好约束下的新形态。
\begin{itemize}
\item 资料版本限制（version risk）：课程页、公开视频和配套资料可能在学期中更新；引用时应注明访问日期或资料版本。
\item 测量口径限制（measurement risk）：FLOPs、tokens/s、benchmark score、reward 等数字只有在协议完整时可比较。
\item 规模外推限制（scaling risk）：小模型或小 batch 的机制不必然外推到 frontier-scale。
\item 数据分布限制（data risk）：数据来源、过滤和去重策略会改变模型行为，也会改变 evaluation contamination 风险。
\item 对齐风险（alignment risk）：post-training 或 multimodal tuning 的行为变化可能提高 helpfulness，同时降低校准、透明度或安全边界。
\end{itemize}
\section{复习要点}
下面的问题把本章内容压缩为可反复检查的条目。每个条目都应能回到前文的解释、公式、伪代码或例题。
\begin{longtable}{p{0.16\linewidth}p{0.38\linewidth}p{0.38\linewidth}}
\toprule
\textbf{条目} & \textbf{对象} & \textbf{检查动作}\\
\midrule
条目 1 & post-training（后训练） & 定义它；写出一个最小例子；说明一个 failure mode；指出它影响哪个资源或评估账本。\\
条目 2 & supervised fine-tuning (SFT) & 定义它；写出一个最小例子；说明一个 failure mode；指出它影响哪个资源或评估账本。\\
条目 3 & RLHF & 定义它；写出一个最小例子；说明一个 failure mode；指出它影响哪个资源或评估账本。\\
条目 4 & reward model（奖励模型） & 定义它；写出一个最小例子；说明一个 failure mode；指出它影响哪个资源或评估账本。\\
条目 5 & DPO & 定义它；写出一个最小例子；说明一个 failure mode；指出它影响哪个资源或评估账本。\\
条目 6 & PPO & 定义它；写出一个最小例子；说明一个 failure mode；指出它影响哪个资源或评估账本。\\
\bottomrule
\end{longtable}
\subsection{综合自测}
\begin{itemize}
\item 我能否不用课程原句，独立解释本章中心问题？
\item 我能否把本章至少一个公式改写成可运行的估算函数？
\item 我能否指出一个看似改进但实际转移成本的方法？
\item 我能否回到视频和课程资料，找到支持本章关键论断的位置？
\item 我能否设计一个最小 ablation，验证本章一个 caveat？
\end{itemize}
\section{总结与延伸}
\subsection{本章小结}
\begin{itemize}
\item Mid/Post-training：SFT 与 RLHF 的核心不是记忆术语，而是理解对象、变量和机制之间的关系。
\item 视频给出讲解顺序，课程资料提供可核对的公式、图、代码或页面结构。
\item 复习时应把每个术语连接到至少一个变量、一个实现动作和一个失败模式。
\end{itemize}
\subsection{延伸解释}
本章中的延伸解释用于连接课程材料中的概念，例如说明公式里的 proxy、解释系统 tradeoff，或把多个材料片段串成一个完整推理链。若需要严肃引用，应回到课程视频和配套资料核对原始表述。
\subsection{练习}
\begin{itemize}
\item 定义题：用中英双语解释本章任意五个重要术语，并说明它们属于 compute、memory、data、optimization、evaluation 还是 deployment 账本。
\item 推导题：选择本章一个公式，逐个解释符号、单位和近似假设，并说明哪个变量最难在真实系统中测量。
\item 实现题：把本章伪代码改写成一个最小 Python sanity check，不要求训练大模型，但要输出可检查的中间量。
\item 资料题：从“课程资料与回看路径”中选择一个视频时间段，回到原始视频核对讲者例子，并记录是否存在字幕误差。
\item 评估题：为本章方法设计一个 ablation 或 benchmark，说明控制变量、数据版本、指标和 failure criteria。
\item 边界题：说明本章一个结论在更长 context、更大 batch、更多 GPU 或更复杂数据分布下可能如何失效。
\item 综合题：把本章内容和前一章或后一章连接起来，写出一个端到端训练/推理 pipeline 中的依赖关系。
\item 批判题：指出一个容易被营销材料夸大的说法，并用本章的资源账本或评估协议反驳它。
\end{itemize}
\end{document}
